% TOP_PROBLEM: {"id": 30007170, "title": "Singularity of Bernoulli convolutions at quartic Salem reciprocals", "category_id": 11, "category": {"id": 11, "name": "dynamical_systems", "display_name": "Dynamical Systems", "description": "Problems about long-term behavior of deterministic systems, Hamiltonian dynamics, and periodic orbits.", "slug": "dynamical-systems", "order_index": 11, "created_at": "2026-07-31T15:26:25.671Z"}, "status": "open"}
% FIRST_POSTED: 2026-10-06
% ENTRY_KIND: statement_only
% !TeX program = lualatex
\documentclass[11pt]{article}
\usepackage[margin=25mm]{geometry}
\usepackage{fontspec}
\setmainfont{Latin Modern Roman}
\usepackage{amsmath,amssymb,mathtools}
\usepackage{unicode-math}
\setmathfont{Latin Modern Math}
\usepackage{xurl,hyperref}
\hypersetup{colorlinks=true,urlcolor=blue}
\setcounter{secnumdepth}{0}
\setlength{\parindent}{0pt}
\setlength{\parskip}{5pt}
\setlength{\emergencystretch}{3em}
\newcommand{\R}{\mathbb R}
\newcommand{\N}{\mathbb N}
\newcommand{\Z}{\mathbb Z}
\newcommand{\Q}{\mathbb Q}
\newcommand{\C}{\mathbb C}
\newcommand{\D}{\mathbb D}
\newcommand{\F}{\mathbb F}
\newcommand{\sref}[2]{\hyperlink{source:#1:#2}{[S#2]}}
\newcommand{\cataloguescope}[1]{#1}
\begin{document}
\section*{Singularity of Bernoulli convolutions at quartic Salem reciprocals}\label{30007170}
\subsection{Definitions and mathematical statement}
\paragraph{Statement recorded in the supplied source.}
For every quartic Salem algebraic integer beta in (1,2), is the Bernoulli convolution given by the distribution of $sum_{n>=0}$ $\epsilon _n$ $\beta ^{-n}$, with independent equiprobable $\epsilon _n$ in \{-1,1\}, singular with respect to Lebesgue measure?
\subsection{Short English statement}
For every quartic Salem algebraic integer beta in (1,2), is the Bernoulli convolution given by the distribution of $sum_{n>=0}$ $\epsilon _n$ $\beta ^{-n}$, with independent equiprobable $\epsilon _n$ in \{-1,1\}, singular with respect to Lebesgue measure?
\subsection{Sources}
\begin{itemize}
\item[S1] Original question source or mathematical context.
\url{https://doi.org/10.4171/owr/2016/21}.
\item[S2] Original question source or mathematical context.
\url{https://doi.org/10.4171/owr/2017/47}.
\item[S3] Original question source or mathematical context.
\url{https://londmathsoc.onlinelibrary.wiley.com/doi/full/10.1112/jlms.70515}.
\item[S4] OpenAI, manuscript citation page. Source of the formulation motivating this local problem.
\url{https://github.com/openai/math/blob/adc7f1241b42e322a6451854ab7e4b4c146bf78a/preprints/Arithmetic-classification-and-non-Pisot-singularity-for-Bernoulli-convolutions-October-3-2026/README.md}.
\end{itemize}
\end{document}
